% Raccolta di esercizi — capitolo 4
% Il capitolo è determinato dal nome di questo file.

\begin{exo}[Integrare una forma differenziale lungo un cammino]{integrale-forme-differentielle}
\keywords{forma differenziale, differenziale esatto, criterio di Schwarz, integrale curvilineo}

Si considerino le due forme differenziali
\[\omega_1=y\,\mathrm{d}x+x\,\mathrm{d}y,\qquad\omega_2=y\,\mathrm{d}x,\]
e tre cammini che congiungono l'origine $O(0,0)$ al punto $M(1,1)$: $\gamma_1$,
formato dal segmento orizzontale $O\to(1,0)$ seguito da quello verticale
$(1,0)\to M$; $\gamma_2$, formato dal segmento verticale $O\to(0,1)$ seguito
da quello orizzontale $(0,1)\to M$; e $\gamma_3$, il segmento rettilineo che
congiunge direttamente $O$ a $M$.

\begin{enumerate}
\item Mostrare che $\omega_1$ è esatta trovando una funzione $f$ tale che $\omega_1=\mathrm{d}f$.
\item Calcolare $\int_{\gamma_1}\omega_1$ e $\int_{\gamma_2}\omega_1$ parametrizzando ogni segmento. Ritrovare il risultato senza calcoli.
\item Applicare il criterio di Schwarz a $\omega_2$. Che cosa si può concludere?
\item Calcolare $\int_{\gamma_1}\omega_2$, $\int_{\gamma_2}\omega_2$ e $\int_{\gamma_3}\omega_2$. Commentare.
\end{enumerate}

\begin{indication}
Su un segmento orizzontale $\mathrm{d}y=0$, su uno verticale $\mathrm{d}x=0$:
ogni integrale si riduce a un integrale semplice.
\end{indication}

\begin{solution}
\textbf{Domanda 1.} Va bene $f(x,y)=xy$, poiché $\partial f/\partial x=y$ e $\partial f/\partial y=x$.

\textbf{Domanda 2.} Per una curva parametrizzata da $t\mapsto\bigl(x(t),y(t)\bigr)$,
\[\int_\gamma\omega_1=\int\left[y(t)x'(t)+x(t)y'(t)\right]\,\mathrm{d}t.\]
Il cammino $\gamma_1$ è l'unione di due segmenti. Sul primo, $x(t)=t$ e
$y(t)=0$, con $t\in[0,1]$; quindi $\mathrm{d}x=\mathrm{d}t$ e $\mathrm{d}y=0$.
Sul secondo, $x(t)=1$ e $y(t)=t$, sempre con $t\in[0,1]$; quindi
$\mathrm{d}x=0$ e $\mathrm{d}y=\mathrm{d}t$. Pertanto,
\[\begin{aligned}I_1=\int_{\gamma_1}\omega_1&=\int_0^1(0\times1+t\times0)\,\mathrm{d}t+\int_0^1(t\times0+1\times1)\,\mathrm{d}t\\&=0+\int_0^1\mathrm{d}t=1.\end{aligned}\]
Per $\gamma_2$, il primo segmento è parametrizzato da $x(t)=0$, $y(t)=t$,
e il secondo da $x(t)=t$, $y(t)=1$, con $t\in[0,1]$ in entrambi i casi. Si ottiene
\[\begin{aligned}I_2=\int_{\gamma_2}\omega_1&=\int_0^1(t\times0+0\times1)\,\mathrm{d}t+\int_0^1(1\times1+t\times0)\,\mathrm{d}t\\&=0+\int_0^1\mathrm{d}t=1.\end{aligned}\]
I due valori coincidono. Lo si poteva prevedere senza calcoli: poiché
$\omega_1=\mathrm{d}f$ con $f(x,y)=xy$, l'integrale dipende solo dagli estremi,
\[\int_\gamma\omega_1=f(M)-f(O)=f(1,1)-f(0,0)=1.\]

\textbf{Domanda 3.} Scriviamo $\omega_2=A(x,y)\,\mathrm{d}x+B(x,y)\,\mathrm{d}y$.
Qui $A(x,y)=y$ e $B(x,y)=0$. Se la forma fosse esatta, il criterio di Schwarz darebbe
\[\frac{\partial A}{\partial y}=\frac{\partial B}{\partial x}.\]
Ma $\partial A/\partial y=1$, mentre $\partial B/\partial x=0$. Le derivate
incrociate sono diverse: $\omega_2$ non è dunque esatta. Il suo integrale può
dipendere dal cammino seguito tra $O$ e $M$, come verifica il calcolo seguente.

\textbf{Domanda 4.} Poiché $\omega_2=y\,\mathrm{d}x$, solo una variazione di
$x$ effettuata a un valore non nullo di $y$ può contribuire all'integrale.
Riprendendo le parametrizzazioni precedenti, su $\gamma_1$ si trova
\[\begin{aligned}J_1=\int_{\gamma_1}\omega_2&=\int_0^1 0\times1\,\mathrm{d}t+\int_0^1 t\times0\,\mathrm{d}t\\&=0.\end{aligned}\]
Su $\gamma_2$, il primo segmento soddisfa $x(t)=0$, dunque $\mathrm{d}x=0$;
sul secondo, $x(t)=t$, $y(t)=1$ e $\mathrm{d}x=\mathrm{d}t$. Quindi
\[\begin{aligned}J_2=\int_{\gamma_2}\omega_2&=\int_0^1 t\times0\,\mathrm{d}t+\int_0^1 1\times1\,\mathrm{d}t\\&=1.\end{aligned}\]
Infine, il segmento diagonale $\gamma_3$ è parametrizzato esplicitamente da
\[x(t)=t,\qquad y(t)=t,\qquad t\in[0,1],\]
da cui $\mathrm{d}x=\mathrm{d}t$ e $\mathrm{d}y=\mathrm{d}t$. Ne segue
\[J_3=\int_{\gamma_3}\omega_2=\int_0^1 y(t)x'(t)\,\mathrm{d}t=\int_0^1 t\,\mathrm{d}t=\frac12.\]
I tre cammini hanno gli stessi estremi, ma danno tre valori diversi: $J_1=0$,
$J_2=1$ e $J_3=1/2$. È precisamente la dipendenza dal cammino caratteristica
di una forma differenziale non esatta.
\end{solution}
\end{exo}

\begin{exo}[Stessa variazione di energia interna, tre modalità di trasferimento]{meme-delta-u-trois-transferts}
\keywords{primo principio, energia interna, calore, lavoro meccanico, lavoro elettrico, calorimetria, funzione di stato}

Un liquido, il suo calorimetro e una piccola resistenza immersa costituiscono
un sistema chiuso, macroscopicamente in quiete. Il recipiente è rigido e si
suppone costante la capacità termica totale del sistema:
\[C=2{,}00\ \mathrm{kJ\,K^{-1}}.\]
Si vuole portare il sistema dallo stato di equilibrio $A$, a
$T_A=20{,}0\,^\circ\mathrm{C}$, allo stato di equilibrio $B$, a
$T_B=25{,}0\,^\circ\mathrm{C}$. Si ammette che
\[U(B)-U(A)=C(T_B-T_A).\]
Le variazioni di energia cinetica e potenziale macroscopiche sono nulle. Si
trascurano tutte le perdite verso l'esterno.

Si realizzano indipendentemente i tre protocolli seguenti, a partire dallo
stesso stato iniziale $A$:
\begin{itemize}
\item[Protocollo 1] Il sistema è posto a contatto termico con una sorgente calda, senza che gli venga fornito alcun lavoro.
\item[Protocollo 2] Le pareti sono termicamente isolate. Una paletta agita il liquido grazie alla caduta di una massa $m$ da un'altezza $h=5{,}00$ m. La paletta e il liquido terminano in quiete, e tutta l'energia meccanica persa dalla massa è trasferita al sistema. Si assuma $g=9{,}81\ \mathrm{m\,s^{-2}}$.
\item[Protocollo 3] Il sistema riceve un calore $Q_3=4{,}00$ kJ e l'energia restante viene fornita elettricamente alla resistenza. La potenza elettrica ricevuta è costante e vale $\mathcal P=300$ W.
\end{itemize}

\begin{enumerate}
\item Calcolare la variazione di energia interna $\Delta U=U(B)-U(A)$ comune ai tre protocolli.
\item Per ciascuno dei primi due protocolli, determinare il calore $Q$ e il lavoro $W$ ricevuti dal sistema. Nel secondo, calcolare la massa $m$ necessaria.
\item Per il terzo protocollo, calcolare il lavoro elettrico ricevuto e la durata dell'alimentazione della resistenza.
\item Le tre trasformazioni hanno gli stessi stati iniziale e finale. Spiegare perché i loro valori di $Q$ e $W$ possono tuttavia essere diversi. Il sistema contiene «calore» o «lavoro» nello stato finale $B$?
\end{enumerate}

\begin{indication}
Applicare $\Delta U=Q+W$ con la convenzione del banchiere. Per la paletta, il
lavoro ricevuto è uguale alla diminuzione $mgh$ dell'energia potenziale della
massa. L'energia elettrica che attraversa il confine attraverso i fili è
conteggiata come lavoro elettrico.
\end{indication}

\begin{solution}
\textbf{Domanda 1.} La differenza di temperatura vale $T_B-T_A=5{,}0$ K. Pertanto,
\[\Delta U=C(T_B-T_A)=2{,}00\times5{,}0=10{,}0\ \mathrm{kJ}.\]
Questo valore dipende soltanto dai due stati di equilibrio $A$ e $B$.

\textbf{Domanda 2.} Nel primo protocollo non si riceve lavoro: $W_1=0$.
Il primo principio dà quindi
\[Q_1=\Delta U=10{,}0\ \mathrm{kJ}.\]
Il calore è positivo perché l'energia entra nel sistema.

Nel secondo protocollo le pareti sono termicamente isolate, quindi $Q_2=0$.
L'intera variazione di energia interna deriva dal lavoro meccanico della paletta:
\[W_2=\Delta U=10{,}0\ \mathrm{kJ}.\]
Poiché $W_2=mgh$,
\[m=\frac{W_2}{gh}=\frac{10{,}0\times10^3}{9{,}81\times5{,}00}\simeq204\ \mathrm{kg}.\]
L'elevato ordine di grandezza ricorda che anche un modesto aumento di
temperatura corrisponde già a un'importante energia meccanica.

\textbf{Domanda 3.} Il primo principio impone
\[W_3=\Delta U-Q_3=10{,}0-4{,}00=6{,}00\ \mathrm{kJ}.\]
Questo lavoro è positivo: l'alimentazione elettrica fornisce energia al sistema.
Con $W_3=\mathcal P\,\Delta t$,
\[\Delta t=\frac{W_3}{\mathcal P}=\frac{6{,}00\times10^3}{300}=20{,}0\ \mathrm{s}.\]

\textbf{Domanda 4.} I tre cammini danno rispettivamente
\[(Q_1,W_1)=(10{,}0,0)\ \mathrm{kJ},\qquad(Q_2,W_2)=(0,10{,}0)\ \mathrm{kJ},\]
e
\[(Q_3,W_3)=(4{,}00,6{,}00)\ \mathrm{kJ}.\]
In tutti i casi, la somma $Q+W$ vale $10{,}0$ kJ e produce la stessa variazione
$\Delta U$. L'energia interna è una funzione di stato, mentre calore e lavoro
descrivono trasferimenti durante una trasformazione. Nello stato finale $B$,
il sistema possiede energia interna, ma non contiene né «calore» né «lavoro».
\end{solution}
\end{exo}

\begin{exo}[Compressione a pressione esterna costante]{compression-pression-exterieure-constante}
\seoready{true}
\keywords{lavoro delle forze di pressione, pressione esterna costante, compressione, espansione, espansione libera}

Un gas viene compresso da un volume $V_A$ a un volume $V_B<V_A$ sotto una
pressione esterna costante $P_0$.

\begin{enumerate}
\item Calcolare il lavoro ricevuto dal gas.
\item Il gas ritorna da $V_B$ a $V_A$ contro la stessa pressione esterna $P_0$. Calcolare il lavoro ricevuto e confrontarlo con quello della compressione.
\item Ripetere l'espansione da $V_B$ a $V_A$ quando il gas si espande nel vuoto. Interpretare i tre segni ottenuti.
\end{enumerate}

\begin{indication}
Usare $W=-\int P_{\mathrm{ext}}\,\mathrm{d}V$. Occorre integrare la pressione
\emph{esterna}, anche quando la trasformazione non è quasistatica.
\end{indication}

\begin{solution}
\textbf{Domanda 1.} La pressione esterna è costante, quindi
\[W_{A\to B}=-P_0(V_B-V_A)=P_0(V_A-V_B)>0.\]
Il gas riceve lavoro durante la compressione.

\textbf{Domanda 2.} Per il ritorno,
\[W_{B\to A}=-P_0(V_A-V_B)<0.\]
I due lavori sono opposti perché le trasformazioni avvengono contro la stessa
pressione esterna costante.

\textbf{Domanda 3.} Nel vuoto, $P_{\mathrm{ext}}=0$, quindi
\[W_{B\to A}^{\mathrm{vide}}=0.\]
Un'espansione non implica dunque automaticamente un lavoro negativo: occorre
che una forza esterna venga effettivamente respinta.
\end{solution}
\end{exo}
